% =====================================================================
% CAP. 11 - NIVEL 2 - QUESTOES GRAFICAS E INTEGRADAS
% Inspiradas no nivel de variedade de livros-texto, com enunciados originais.
% =====================================================================
\blocoaplicacao

\begin{questao}[2]{}
A chave do circuito permaneceu fechada por muito tempo e \textbf{abre} em
$t=0$. Determine $v_C(0^+)$, a constante de tempo para $t>0$ e
$v_C(\SI{0.30}{\second})$.

\circuitoqq{%
\begin{circuitikz}[scale=0.96,transform shape,line width=0.8pt,every node/.append style={font=\scriptsize}]
\draw (0,0) to[american voltage source,l_={$\SI{18}{\volt}$}] (0,3)
      to[switch] (1.9,3)
      to[R,l={$R_1=\SI{3}{\kilo\ohm}$}] (4.5,3) coordinate(a);
\node[anchor=south] at (0.95,3.30) {abre em $t=0$};
\draw (a) to[R,l_={$R_2=\SI{6}{\kilo\ohm}$}] (4.5,0) -- (0,0);
\draw (a) -- (7.0,3)
      to[C,v^>={$v_C$}] (7.0,0) -- (4.5,0);
\node[anchor=east] at (6.75,1.5) {$C=\SI{100}{\micro\farad}$};
\end{circuitikz}}{Circuito RC com mudança de topologia em $t=0$.}

\resposta{$v_C(0^+)=\SI{12}{\volt}$; $\tau=\SI{0.60}{\second}$;
$v_C(0{,}30\,\mathrm{s})\approx\SI{7.28}{\volt}$}
\resolucao{Para $t<0$, o capacitor está em regime permanente e equivale a um
circuito aberto. Assim,
\[
v_C(0^-)=18\frac{6}{3+6}=\SI{12}{\volt}.
\]
Pela continuidade da tensão no capacitor, $v_C(0^+)=\SI{12}{\volt}$. Quando a
chave abre, a fonte e $R_1$ ficam desconectados do capacitor; a descarga ocorre
somente por $R_2$. Logo
\[
\tau=R_2C=(6\times10^3)(100\times10^{-6})=\SI{0.60}{\second}.
\]
Como $v_C(\infty)=0$,
$v_C(t)=12e^{-t/0{,}60}$ e, em $t=\SI{0.30}{\second}$,
$v_C\approx12e^{-0{,}5}=\SI{7.28}{\volt}$.}
\end{questao}

\begin{questao}[2]{}
No circuito abaixo, a chave permaneceu fechada por muito tempo e abre em
$t=0$. Determine $i_L(0^+)$, a constante de tempo da descarga e
$i_L(\SI{10}{\milli\second})$.

\circuitoqq{%
\begin{circuitikz}[scale=0.96,transform shape,line width=0.8pt,every node/.append style={font=\scriptsize}]
\draw (0,0) to[american voltage source,l_={$\SI{24}{\volt}$}] (0,3)
      to[switch] (1.9,3)
      to[R,l={$R_1=\SI{6}{\ohm}$}] (4.5,3) coordinate(a);
\node[anchor=south] at (0.95,3.30) {abre em $t=0$};
\draw (a) to[L,i>^={$i_L$}] (4.5,0) -- (0,0);
\node[anchor=east] at (4.25,1.5) {$L=\SI{120}{\milli\henry}$};
\draw (a) -- (7.0,3) to[R,l_={$R_2=\SI{18}{\ohm}$}] (7.0,0) -- (4.5,0);
\end{circuitikz}}{Circuito RL com caminho de descarga resistivo.}

\resposta{$i_L(0^+)=\SI{4}{\ampere}$; $\tau\approx\SI{6.67}{\milli\second}$;
$i_L(10\,\mathrm{ms})\approx\SI{0.893}{\ampere}$}
\resolucao{Em regime permanente, o indutor equivale a um curto-circuito; portanto
$R_2$ fica sem tensão e
$i_L(0^-)=24/6=\SI{4}{\ampere}$. Pela continuidade da corrente,
$i_L(0^+)=\SI{4}{\ampere}$. Após a abertura da chave, a fonte e $R_1$ ficam
isolados e o indutor descarrega por $R_2$:
\[
\tau=\frac{L}{R_2}=\frac{0{,}12}{18}\approx\SI{6.67}{\milli\second}.
\]
Assim, $i_L(t)=4e^{-t/\tau}$ e
$i_L(10\,\mathrm{ms})\approx4e^{-10/6{,}67}=\SI{0.893}{\ampere}$.}
\end{questao}

\begin{questao}[2]{}
A fonte de corrente já está ligada há muito tempo. Em $t=0$ a chave fecha,
inserindo $R_2$ em paralelo com o capacitor. Determine $v_C(0^+)$,
$v_C(\infty)$, $\tau$ e escreva $v_C(t)$ para $t>0$.

\circuitoqq{%
\begin{circuitikz}[scale=0.94,transform shape,line width=0.8pt,every node/.append style={font=\scriptsize}]
\draw (0,0) to[american current source,l_={$\SI{2}{\milli\ampere}$}] (0,3) -- (2.2,3) coordinate(a);
\draw (a) to[R,l_={$R_1=\SI{4}{\kilo\ohm}$}] (2.2,0) -- (0,0);
\draw (a) -- (4.7,3) to[C,v^>={$v_C$}] (4.7,0) -- (2.2,0);
\node[anchor=east] at (4.45,1.5) {$C=\SI{50}{\micro\farad}$};
\draw (a) -- (6.0,3) to[switch] (7.5,3)
      to[R,l_={$R_2=\SI{12}{\kilo\ohm}$}] (7.5,0) -- (4.7,0);
\node[anchor=south] at (6.75,3.30) {fecha em $t=0$};
\end{circuitikz}}{Mudança da resistência vista pelo capacitor.}

\resposta{$v_C(0^+)=\SI{8}{\volt}$; $v_C(\infty)=\SI{6}{\volt}$;
$\tau=\SI{150}{\milli\second}$;
$v_C(t)=6+2e^{-t/0{,}15}\ \si{\volt}$}
\resolucao{Antes da comutação, o capacitor está aberto e $R_2$ desconectado. Logo
$v_C(0^-)=I_SR_1=(2\,\mathrm{mA})(4\,\mathrm{k\Omega})=\SI{8}{\volt}$, portanto
$v_C(0^+)=\SI{8}{\volt}$. Após o fechamento, em regime permanente o capacitor
continua aberto e
\[
v_C(\infty)=I_S(R_1\parallel R_2)
=2\,\mathrm{mA}\,(3\,\mathrm{k\Omega})=\SI{6}{\volt}.
\]
Para a constante de tempo, a fonte de corrente independente é aberta:
$R_{\mathrm{eq}}=R_1\parallel R_2=\SI{3}{\kilo\ohm}$ e
$\tau=R_{\mathrm{eq}}C=\SI{0.15}{\second}$. Assim,
$v_C(t)=6+(8-6)e^{-t/0{,}15}$.}
\end{questao}

\begin{questao}[2]{}
A chave fecha em $t=0$, aplicando um degrau de $\SI{20}{\volt}$ ao circuito RLC
em série. Sem resolver ainda a resposta completa, determine $\alpha$, $\omega_0$,
$\omega_d$ e classifique o amortecimento.

\circuitoqq{%
\begin{circuitikz}[scale=0.92,transform shape,line width=0.8pt,every node/.append style={font=\scriptsize}]
\draw (0,0) to[american voltage source,l_={$\SI{20}{\volt}$}] (0,3)
      to[switch] (1.8,3)
      to[R,l={$R=\SI{15}{\ohm}$}] (4.0,3)
      to[L,l={$L=\SI{50}{\milli\henry}$}] (6.5,3)
      to[C] (9.0,3) -- (9.0,0) -- (0,0);
\node[anchor=south] at (0.9,3.30) {fecha em $t=0$};
\node[anchor=south] at (7.75,3.30) {$C=\SI{200}{\micro\farad}$};
\end{circuitikz}}{Circuito RLC em série excitado por degrau.}

\resposta{$\alpha=\SI{150}{\per\second}$;
$\omega_0\approx\SI{316.2}{\radian\per\second}$;
$\omega_d\approx\SI{278.4}{\radian\per\second}$; resposta subamortecida}
\resolucao{Para RLC em série,
\[
\alpha=\frac{R}{2L}=\frac{15}{0{,}10}=\SI{150}{\per\second},\qquad
\omega_0=\frac{1}{\sqrt{LC}}
=\frac{1}{\sqrt{0{,}05\cdot200\times10^{-6}}}
\approx\SI{316.2}{\radian\per\second}.
\]
Como $\alpha<\omega_0$, a resposta é subamortecida e
$\omega_d=\sqrt{\omega_0^2-\alpha^2}\approx\SI{278.4}{\radian\per\second}$.}
\end{questao}

\begin{questao}[2]{}
O circuito RLC em paralelo abaixo parte sem energia armazenada e recebe a fonte
$i_s(t)=3u(t)\,\si{\ampere}$. Determine $v(0^+)$, $i_L(0^+)$,
$\mathrm{d}v/\mathrm{d}t|_{0^+}$, $v(\infty)$ e $i_L(\infty)$.

\circuitoqq{%
\begin{circuitikz}[scale=0.96,transform shape,line width=0.8pt,every node/.append style={font=\scriptsize}]
\draw (0,0) to[american current source] (0,3) -- (7.2,3);
\node[anchor=east] at (-0.25,1.5) {$3u(t)\,\si{\ampere}$};
\draw (2.2,3) to[R] (2.2,0);
\node[anchor=east] at (1.95,1.5) {$R=\SI{6}{\ohm}$};
\draw (4.7,3) to[C,v^>={$v$}] (4.7,0);
\node[anchor=east] at (4.45,1.5) {$C=\SI{0.5}{\farad}$};
\draw (7.2,3) to[L,i>^={$i_L$}] (7.2,0) -- (0,0);
\node[anchor=east] at (6.95,1.5) {$L=\SI{2}{\henry}$};
\end{circuitikz}}{RLC em paralelo: condições iniciais, derivada inicial e regime final.}

\resposta{$v(0^+)=0$; $i_L(0^+)=0$;
$\dot v(0^+)=\SI{6}{\volt\per\second}$;
$v(\infty)=0$; $i_L(\infty)=\SI{3}{\ampere}$}
\resolucao{Como o circuito parte sem energia, $v_C(0^-)=0$ e $i_L(0^-)=0$;
logo $v(0^+)=0$ e $i_L(0^+)=0$. Aplicando KCL em $0^+$,
\[
3=\frac{v}{R}+C\dot v+i_L
\quad\Rightarrow\quad
3=0+0{,}5\dot v+0,
\]
portanto $\dot v(0^+)=\SI{6}{\volt\per\second}$. Em regime CC, o capacitor é
aberto e o indutor é curto-circuito; assim $v(\infty)=0$ e toda a corrente da fonte
flui pelo indutor: $i_L(\infty)=\SI{3}{\ampere}$.}
\end{questao}

\gabarito[Nível 2 --- circuitos e chaveamento]
